\documentclass[12pt]{article}

\usepackage[utf8]{inputenc}
\usepackage{geometry}
\geometry{a4paper,scale=0.75}
\linespread{1.5}
\usepackage{graphicx} 
\usepackage{float} 
\usepackage{subcaption} 
\usepackage{enumerate}
\usepackage{enumitem}
\usepackage{amsmath}
\usepackage{array}
\usepackage{booktabs}
\usepackage{multirow}
\usepackage{amsfonts}
\usepackage[english]{babel}
\usepackage{amsthm}
\usepackage{dcolumn}
\usepackage{multicol}
\usepackage{stfloats}
\usepackage{lscape}
\usepackage[figuresright]{rotating}
\RequirePackage{pdflscape}
\usepackage[toc,page]{appendix}
\usepackage{geometry}
\usepackage{longtable}
\usepackage{comment}
\usepackage{xcolor}

% -------- enumerated sub-labels (a), (b), … --
\usepackage{enumitem}
\setlist[enumerate,1]{label=(\alph*),ref=\alph*}
% ---------------------------------------------

\usepackage{hyperref}
\hypersetup{hidelinks,
	colorlinks=true,
	allcolors=black,
	pdfstartview=Fit,
	breaklinks=true}
\usepackage{csquotes}
\usepackage{natbib}
\bibliographystyle{apalike}
\newtheorem{definition}{Definition}
\newtheorem{theorem}{Theorem}
\newtheorem{proposition}[theorem]{Proposition}
\newtheorem{lemma}[theorem]{Lemma}
\newtheorem{corollary}[theorem]{Corollary}
\newtheorem*{remark}{Remark}
\newtheorem{example}{Example}
\newtheorem{exercise}{Exercise}
\newtheorem{assumption}{Assumption}[section] % number within sections

\begin{document}

\begin{center}
    ECON 5260: Macroeconomic Theory II\\
    {\large \textbf{Grading Criteria for the Literature Review}}\\
    Teaching Assistant: Harlly Zhou
\end{center}


The literature review is expected to be 10--15 pages long. The full mark for each item below is 5 points. The final grade will be calculated according to the following weights: Content accounts for 50\%, Structure accounts for 30\%, and Writing accounts for 20\%.

\subsection*{1. Content \texorpdfstring{(50\%)}{}}

This criterion evaluates the substantive quality of the literature review. A high-quality review should cover a sufficient range of relevant papers and demonstrate a clear understanding of the research area. In particular, the review should include the classical or foundational references that establish the field, while also discussing recent progress in the literature.

The review should accurately capture the main contribution, model, mechanism, and findings of each selected paper. It should not simply summarize papers one by one at a superficial level. Instead, it should show that the student has internalized the models and arguments in the literature. A strong review explains the economic intuition behind the papers, identifies how different papers are connected, and discusses how the literature has developed over time.

Furthermore, a strong literature review should not only summarize the selected papers, but also evaluate them. It should discuss the limitations of the existing literature, compare the assumptions and mechanisms across papers, and identify unresolved questions. A review that only describes papers one by one, without critique or synthesis, should receive a lower score in both Content and Structure.

A lower score should be assigned if the review omits important references, relies on too few papers, focuses only on recent or only on classical papers, misrepresents the papers discussed, or appears to copy formulas, model setups, or descriptions from the literature without demonstrating understanding.

\subsection*{2. Structure \texorpdfstring{(30\%)}{}}

This criterion evaluates the organization and coherence of the literature review. A complete literature review should contain the standard components of an academic review, including an abstract, an introduction, the main text, a conclusion, and a reference list.

The introduction should clearly define the topic, explain why the topic is important, and preview the organization of the review. The main text should be logically structured around themes, mechanisms, models, assumptions, empirical or theoretical approaches, or contributions, rather than simply listing papers in sequence. A strong review synthesizes the literature by comparing different specifications, identifying common frameworks, explaining disagreements or extensions, and showing how each paper contributes to the broader research question.

The conclusion should summarize the main lessons from the literature and, when appropriate, point out open questions or possible directions for future research. A lower score should be assigned if the review lacks key sections, has weak transitions, piles up papers without synthesis, or fails to organize the literature around a clear logic.

\subsection*{3. Writing \texorpdfstring{(20\%)}{}}

This criterion evaluates the academic quality of the writing. The review should follow an appropriate academic format, with clear section headings, consistent citation style, complete references, and proper use of equations, tables, or figures when needed.

The language should be accurate, clear, and academic. The writing should avoid informal expressions, vague claims, unsupported statements, and excessive repetition. The logic of the writing should be easy to follow, with each paragraph serving a clear purpose and each claim supported by relevant discussion of the literature.

A lower score should be assigned if the review contains frequent grammar problems, unclear sentences, inconsistent formatting, inaccurate terminology, weak paragraph structure, or citations and references that are incomplete or improperly formatted.

\end{document}